\documentclass{situationist-preprint}

\title{The Size of the Universe of Numbers}
\msid{MTH-01-c}
\msfamily{MTH-01}
\msdiscipline{Philosophy of Mathematics}
\mstype{Consequence}
\mskeywords{indispensability; real numbers; cardinality; continuum; transfinite numbers; Burali-Forti}

\begin{document}
\maketitle

\begin{abstract}
If, as MTH-01 holds, numbers are abstract objects whose existence the indispensability argument establishes, how many numbers are there? We derive three consequences. First, the argument establishes a lower bound: the real and complex numbers used in physics exist as completed totalities, so there are at least continuum many numbers, and, given MTH-06, more than $\aleph_1$. Second, the argument establishes no upper bound, and the minimalist reading on which only the numbers physics uses exist is unstable. Third, once the transfinite numbers of set theory are admitted, the numbers form no set, so that the universe of numbers has no size at all.
\end{abstract}

\section{The lower bound}\label{s:lower}

The indispensability argument establishes existence by establishing the truth of the theories in which numbers are used. Its conclusion is therefore as large as the mathematics of those theories. Classical mechanics represents positions and velocities by real numbers; electromagnetism uses real-valued fields on space; quantum mechanics represents states by vectors in complex Hilbert spaces, and the values of observables by real numbers. The commitment is not to the numbers that happen to be the values of measured magnitudes, which are finitely many, but to the domain over which the theories quantify: the laws state that for every time $t$ in an interval of reals, the state satisfies a differential equation. MTH-01 accordingly yields, as MTH-10 uses it, the existence of the completed infinite sets used in science, and among them the field of real numbers.

\begin{proposition}\label{p:lower}
Given MTH-01, there are at least $2^{\aleph_0}$ numbers. Given MTH-06 in addition, there are more than $\aleph_1$.
\end{proposition}

\begin{proof}
By MTH-01 and the use of the real field in physics, the real numbers exist; there are $2^{\aleph_0}$ of them, uncountably many, as \citet{Cantor1874} showed. By MTH-06, the continuum hypothesis is false, so $2^{\aleph_0} > \aleph_1$.
\end{proof}

The bound is sharp for the number systems used in physics. The naturals, integers, rationals, reals and complex numbers together have cardinality $2^{\aleph_0}$, since each is contained in or constructed from the reals without increase of cardinality. The functions over which physics also quantifies, the fields and the elements of function spaces, are more numerous, but they are not numbers; they are mathematical objects whose existence the argument also supports, and which MTH-01's conclusion does not concern.

\section{No upper bound}\label{s:upper}

\citet{Quine1986} drew a minimalist moral from the argument: mathematics is justified so far as it is applied, together with what is needed to round it out; the higher reaches of set theory, beyond any such need, are recreation without ontological rights. On this reading the universe of numbers would be exactly the numbers of applied mathematics, and so of cardinality exactly $2^{\aleph_0}$.

\begin{proposition}\label{p:noupper}
The indispensability argument does not support the claim that there are no numbers other than those used in physics.
\end{proposition}

\begin{proof}
Principle P1 of MTH-01-b gives a sufficient condition for justified belief in existence, not a necessary one. That a kind of entity is not indispensable to physics is no evidence that it does not exist, any more than the absence of a species from the theories of geology is evidence that it does not exist. A negative existential about numbers would need its own support, and the argument supplies none.
\end{proof}

The minimalist position is also unstable on its own terms. Quine's \enquote{rounding out} admits whatever mathematics is needed to give the applied part a natural closure; the real field is so admitted, since no physical theory uses all reals. But the same consideration admits the countable ordinals: the analysis of closed sets of reals by transfinite iteration of the derived set, a part of real analysis, is the context in which Cantor introduced them \citep{Cantor1883}. Rounding out has no principled stopping point short of the mathematics in which applied mathematics is naturally set, and that is set theory with its transfinite numbers.

\section{No size}\label{s:nosize}

Cantor called the ordinals and cardinals of set theory numbers, and mathematical usage continues to do so. If they are numbers, the question how many numbers there are changes its character.

\begin{proposition}\label{p:nosize}
If the transfinite ordinals are numbers, the numbers do not form a set, and there is no cardinal number of numbers.
\end{proposition}

\begin{proof}
Suppose the ordinals formed a set. It would be well-ordered by membership, and its order type would be an ordinal greater than every ordinal, including itself. This is the Burali-Forti paradox \citep{BuraliForti1897}. So the ordinals form no set, and a fortiori neither do the numbers that include them; a cardinality is assigned only to sets.
\end{proof}

The three results fit together. What the indispensability argument establishes is a lower bound, more than $\aleph_1$ given MTH-06. What it leaves open, set theory fills: once transfinite numbers are admitted, there is no upper bound because there is no size. The universe of numbers is not a large set but an open-ended totality, of which every set of numbers is a proper part. None of this conflicts with MTH-14's thesis that numbers are positions in structures: each number system is a structure, and the totality of positions in all of them is what has no size.

\begin{objection}[Overreach]
MTH-01 concerns the numbers established by indispensability. Proposition~\ref{p:nosize} concerns numbers that physics does not use, and so goes beyond the family's conclusion.
\end{objection}

\begin{reply}
Proposition~\ref{p:nosize} is conditional. It asserts nothing about whether the transfinite ordinals exist; it says what follows for the size of the universe of numbers if they do. Proposition~\ref{p:noupper} shows that the indispensability argument leaves the question open, and Section~\ref{s:upper} gives a reason, internal to the indispensabilist's own practice, for answering it affirmatively. The family's conclusion fixes the lower bound; the consequence drawn here is that it fixes nothing more.
\end{reply}

\section{Conclusion}

The indispensability argument settles that there are at least continuum many numbers, and, with the failure of the continuum hypothesis, more than $\aleph_1$. It sets no upper limit, and the minimalist attempt to make the applied numbers exhaustive does not survive its own rounding-out. If the transfinite numbers are numbers, there is no number of numbers.

\begin{thebibliography}{9}\small
\bibitem[Burali-Forti(1897)]{BuraliForti1897} Burali-Forti, C. (1897). Una questione sui numeri transfiniti. \emph{Rendiconti del Circolo Matematico di Palermo}, 11, 154--164.
\bibitem[Cantor(1874)]{Cantor1874} Cantor, G. (1874). Ueber eine Eigenschaft des Inbegriffes aller reellen algebraischen Zahlen. \emph{Journal f\"ur die reine und angewandte Mathematik}, 77, 258--262.
\bibitem[Cantor(1883)]{Cantor1883} Cantor, G. (1883). \emph{Grundlagen einer allgemeinen Mannichfaltigkeitslehre}. Teubner.
\bibitem[Quine(1986)]{Quine1986} Quine, W.~V. (1986). Reply to Charles Parsons. In L.~E. Hahn and P.~A. Schilpp (eds.), \emph{The Philosophy of W.~V. Quine}, 396--403. Open Court.
\end{thebibliography}

\end{document}
